% Original analytic teaching diagram, 2026-11-12.
% n=3 independent trials with success probability p=0.2. The 8 outcome sequences
% are grouped by the number k of successes (filled square = success).
% Per sequence: p^k(1-p)^(3-k) = 0.512, 0.128, 0.032, 0.008;
% column totals C(3,k)p^k(1-p)^(3-k) = 0.512, 0.384, 0.096, 0.008.
\definecolor{VB}{HTML}{7188A3}\definecolor{VBF}{HTML}{E5EAF0}
\definecolor{VG}{HTML}{829C8B}\definecolor{VGF}{HTML}{E8EEE9}
\definecolor{VY}{HTML}{B4A16F}\definecolor{VYF}{HTML}{F1ECDD}
\definecolor{VR}{HTML}{A7787D}\definecolor{VRF}{HTML}{F0E5E6}
\definecolor{VP}{HTML}{9588A5}\definecolor{VPF}{HTML}{ECE8F0}
\definecolor{VI}{HTML}{3B4252}\definecolor{VM}{HTML}{6B7280}\definecolor{VL}{HTML}{4C4D4F}
\begin{tikzpicture}[x=1mm,y=1mm,>=Stealth,font=\figurefont\fontsize{8}{10}\selectfont,
lab/.style={text=VM,font=\figurefont\fontsize{7.5}{9}\selectfont},
s/.style={draw=VB,fill=VB,line width=.6pt},
f/.style={draw=VL,fill=white,line width=.6pt}]
% one sequence: #1 x-left, #2 y-bottom, #3 #4 #5 = 1 (success) or 0
\newcommand{\seqrow}[5]{%
  \foreach \b [count=\q] in {#3,#4,#5}{
    \ifnum\b=1 \draw[s] ({#1+(\q-1)*5},#2) rectangle ++(4,4);
    \else \draw[f] ({#1+(\q-1)*5},#2) rectangle ++(4,4); \fi}}
% column k = 0,1,2,3 at x = 20, 45, 70, 95 (left edge of first square)
\foreach \k/\x in {0/20,1/45,2/70,3/95}{\node[text=VI] at ({\x+7},40) {$k=\k$};}
\seqrow{20}{30}{0}{0}{0}
\seqrow{45}{30}{1}{0}{0}\seqrow{45}{23}{0}{1}{0}\seqrow{45}{16}{0}{0}{1}
\seqrow{70}{30}{1}{1}{0}\seqrow{70}{23}{1}{0}{1}\seqrow{70}{16}{0}{1}{1}
\seqrow{95}{30}{1}{1}{1}
\draw[VL,line width=.5pt] (16,12) -- (113,12);
\node[lab,anchor=east] at (14,7) {每条};
\node[lab,anchor=east] at (14,0) {合计};
\foreach \x/\v in {20/0.512,45/0.128,70/0.032,95/0.008}{\node[text=VI] at ({\x+7},7) {$\v$};}
\foreach \x/\v in {20/0.512,45/0.384,70/0.096,95/0.008}{\node[text=VB] at ({\x+7},0) {$\v$};}
\end{tikzpicture}
